% ====================================================================
% ФАЙЛ: tasks/01_mechanics/work_power_bank_part1.tex
% ТЕМА: РАБОТА И МОЩНОСТЬ СИЛЫ. КИНЕТИЧЕСКАЯ ЭНЕРГИЯ
% ПАЧКА 1: ВАРИАНТЫ 1–10
% ====================================================================

% === ЗАДАЧА 1 ===
% Тег: #МЕХАНИКА #КИМ_03 #ВАРИАНТ_01
\begin{taskbasic}[code={\label{task:work_v1_n3}}]
\textbf{Задача 1.} При равномерном прямолинейном перемещении саней по горизонтальному участку пути на $20$~м постоянная горизонтально направленная сила тяги совершает работу $240$~Дж. Чему равен модуль силы трения?

\kim{3}
\end{taskbasic}
\answer{12}

% === ЗАДАЧА 2 ===
% Тег: #МЕХАНИКА #КИМ_03 #ВАРИАНТ_02
\begin{taskbasic}[code={\label{task:work_v2_n3}}]
\textbf{Задача 2.} Под действием постоянной горизонтально направленной силы тяги $80$~Н ящик равномерно прямолинейно перемещается по горизонтальному участку пути на расстояние $15$~м. Какую работу совершает при этом сила трения?

\kim{3}
\end{taskbasic}
\answer{-1200}

% === ЗАДАЧА 3 ===
% Тег: #МЕХАНИКА #КИМ_03 #ВАРИАНТ_05
\begin{taskbasic}[code={\label{task:work_v5_n3}}]
\textbf{Задача 3.} Скорость волана при подаче в бадминтоне составила $288$~км/ч. Масса волана $6$~г. Какова кинетическая энергия волана при подаче?

\kim{3}
\end{taskbasic}
\answer{19{,}2}

% === ЗАДАЧА 4 ===
% Тег: #МЕХАНИКА #КИМ_03 #ВАРИАНТ_06
\begin{taskbasic}[code={\label{task:work_v6_n3}}]
\textbf{Задача 4.} Какую скорость при подаче сообщили волейбольному мячу, если его кинетическая энергия в этот момент составила $56$~Дж. Масса мяча $280$~г.

\kim{3}
\end{taskbasic}
\answer{20}

% ====================================================================
% ФАЙЛ: tasks/01_mechanics/work_power_bank_part2.tex
% ТЕМА: РАБОТА И МОЩНОСТЬ СИЛЫ
% ПАЧКА 2: ВАРИАНТЫ 11–20
% ====================================================================

% === ЗАДАЧА 1 ===
% Тег: #МЕХАНИКА #КИМ_03 #ВАРИАНТ_11
\begin{taskbasic}[code={\label{task:work_v11_n3}}]
\textbf{Задача 1.} Шарик массой $200$~г падает с высоты $5$~м на горизонтальную плиту. Чему равна работа силы тяжести, действующей на шарик, на всём пути падения? Сопротивлением воздуха пренебречь.

\kim{3}
\end{taskbasic}
\answer{10}

% === ЗАДАЧА 2 ===
% Тег: #МЕХАНИКА #КИМ_03 #ВАРИАНТ_12
\begin{taskbasic}[code={\label{task:work_v12_n3}}]
\textbf{Задача 2.} Шарик массой $300$~г падает с высоты $5$~м на горизонтальную плиту. Чему равна работа силы тяжести, действующей на шарик, на всём пути падения? Сопротивлением воздуха пренебречь.

\kim{3}
\end{taskbasic}
\answer{15}

% === ЗАДАЧА 3 ===
% Тег: #МЕХАНИКА #КИМ_03 #ВАРИАНТ_17
\begin{taskbasic}[code={\label{task:power_v17_n3}}]
\textbf{Задача 3.} Автомобиль движется по прямой горизонтальной дороге с постоянной скоростью $108$~км/ч. Модуль силы сопротивления движению автомобиля равен $12$~кН. Чему равна мощность, развиваемая двигателем автомобиля? Ответ выразите в кВт.

\kim{3}
\end{taskbasic}
\answer{360}

% === ЗАДАЧА 4 ===
% Тег: #МЕХАНИКА #КИМ_03 #ВАРИАНТ_18
\begin{taskbasic}[code={\label{task:power_v18_n3}}]
\textbf{Задача 4.} Автомобиль движется по прямой горизонтальной дороге с постоянной скоростью $72$~км/ч. Модуль силы сопротивления движению автомобиля равен $6$~кН. Чему равна мощность, развиваемая двигателем автомобиля? Ответ выразите в кВт.

\kim{3}
\end{taskbasic}
\answer{120}
% === ЗАДАЧА X ===
% Тег: #МЕХАНИКА #КИМ_03 #РАБОТА_И_МОЩНОСТЬ
\begin{taskbasic}[code={\label{task:power_inclined_force_n3}}]
\textbf{Задача X.} Брусок массой $5$~кг равномерно перемещают по горизонтальной поверхности со скоростью $1$~м/с, прикладывая к нему постоянную силу $4$~Н, направленную под углом $60^{\circ}$ к горизонту. Чему равна мощность силы $F$?

\nopagebreak\smallskip
\begin{center}
\begin{tikzpicture}[scale=1.2, every node/.style={font=\footnotesize}]
    % Поверхность стола (пол)
    \draw[thick] (-1.0,0) -- (2.5,0);
    \foreach \x in {-1.0,-0.8,-0.6,-0.4,-0.2,0,0.2,0.4,0.6,0.8,1.0,1.2,1.4,1.6,1.8,2.0,2.2,2.4}
        \draw[thin] (\x,0) -- (\x-0.1,-0.1);

    % Брусок
    \draw[thick, fill=brandPrimary!20] (0,0) rectangle (1.2,0.6);
    
    % Линия горизонта для угла
    \draw[thick] (1.2,0.2) -- (2.2,0.2);
    
    % Вектор силы F (под углом 60 градусов)
    \draw[-{Stealth[scale=0.8]}, ultra thick, brandAccent] (1.2,0.2) -- ({1.2 + 1.0*cos(60)}, {0.2 + 1.0*sin(60)}) node[left=2pt] {$\vec{F}$};
    
    % Обозначение угла 60 градусов
    \draw (1.5,0.2) arc (0:60:0.3);
    \node at (1.8,0.4) {$60^{\circ}$};
\end{tikzpicture}
\end{center}

\kim{3}
\end{taskbasic}
\answer{2}

% === ЗАДАЧА X ===
% Тег: #МЕХАНИКА #КИМ_03 #РАБОТА_СИЛЫ
\begin{taskbasic}[code={\label{task:work_inclined_friction_n3}}]
\textbf{Задача X.} Брусок массой $2$~кг, к которому приложена сила $4$~Н, направленная вертикально вверх, равномерно движется вниз по шероховатой наклонной плоскости с углом при основании $30^{\circ}$. Чему равен модуль работы, которую совершит над бруском сила трения при перемещении бруска на $1$~м?

\nopagebreak\smallskip
\begin{center}
\begin{tikzpicture}[scale=1.3, every node/.style={font=\footnotesize}]
    % Наклонная плоскость
    \draw[thick, fill=gray!5] (0,0) -- (4.0,0) -- (4.0,2.31) -- cycle; % 4*tan(30) = 2.31
    
    % Угол наклонной плоскости
    \draw (0.5,0) arc (0:30:0.5);
    \node at (0.8,0.22) {$30^{\circ}$};

    % Брусок на склоне (повернутый)
    \begin{scope}[shift={(1.5,0.866)}, rotate=30]
        \draw[thick, fill=brandPrimary!20] (0,0) rectangle (1.2,0.6);
        % Точка приложения вертикальной силы в центре бруска
        \coordinate (C) at (0.6,0.3);
    \end{scope}
    
    % Сила F (строго вертикально вверх из центра бруска)
    \draw[-{Stealth[scale=0.8]}, ultra thick, brandAccent] (C) -- ++(0,1.0) node[right] {$\vec{F}$};
\end{tikzpicture}
\end{center}

\kim{3}
\end{taskbasic}
\answer{8}